% From idea to article: how a research paper gets made
% Build: `make slides` from the repository root (or `latexmk` in this directory).

\documentclass[aspectratio=169,11pt,t]{beamer}

\usepackage[T1]{fontenc}
\usepackage[utf8]{inputenc}
\usepackage[english]{babel}
\usepackage[sfdefault,scaled=0.95]{FiraSans}
\usepackage[scaled=0.85]{FiraMono}
\usepackage{microtype}
\usepackage{booktabs}
\usepackage{array}
\usepackage{tikz}
\usetikzlibrary{positioning,arrows.meta,calc}

\usetheme{seminar}

\hypersetup{colorlinks=true,urlcolor=semAccent,linkcolor=semInk}

\title[From idea to article]{From Idea to Article:\\ How a Research Paper Gets Made}
\subtitle{A first-timer's guide to doing and writing research}
\author{Michał Woźniak}
\institute{Faculty of Economic Sciences, University of Warsaw}
\date{Diploma seminar, academic year 2026/27}

\begin{document}

\begin{frame}[plain,noframenumbering]
  \titlepage
\end{frame}

\begin{frame}{Agenda}
  \tableofcontents
\end{frame}

% =============================================================================
\section{The process}
% =============================================================================

\begin{frame}{Research is a loop, not a line}
  \begin{center}
  \begin{tikzpicture}[
      node distance=3.5mm,
      stage/.style={draw=semRule,fill=semTint,rounded corners=2pt,minimum height=11mm,
                    text width=17mm,align=center,font=\scriptsize},
      arr/.style={-{Stealth[length=2mm]},semAccent,line width=0.9pt}]
    \node[stage] (a) {\key{Read}};
    \node[stage,right=of a] (b) {\key{Question}\\ and gap};
    \node[stage,right=of b] (c) {\key{Design}\\ data, method};
    \node[stage,right=of c] (d) {\key{Experiment}};
    \node[stage,right=of d] (e) {\key{Write}};
    \node[stage,right=of e] (f) {\key{Review}\\ and revise};
    \foreach \x/\y in {a/b,b/c,c/d,d/e,e/f} \draw[arr] (\x) -- (\y);
    \draw[arr,semMuted,dashed] (f.south) -- ++(0,-6mm) -| (b.south);
    \node[font=\scriptsize,text=semMuted] at ($(d.south)+(0,-8.5mm)$) {results change the question; writing exposes gaps in the design};
  \end{tikzpicture}
  \end{center}

  \medskip
  \begin{itemize}
    \item You will go around the loop several times. That is normal, not a sign of failure.
    \item \key{Reading never stops}; it only changes purpose: first to find a question, later to
          position your result.
    \item \key{Writing starts on day one.} The concept note becomes the introduction; method notes
          become the methods section. Nobody ``writes up'' at the end.
    \item The unit of progress is a \key{small, finished, reviewable piece}: one card on the board.
  \end{itemize}
\end{frame}

% =============================================================================
\section{Finding a question}
% =============================================================================

\begin{frame}{Where research ideas come from}
  \begin{columns}[T,onlytextwidth]
    \column{0.5\textwidth}
    \textbf{From the literature}
    \begin{itemize}
      \item Recent issues of the \key{top journals}: which questions are open right now?
      \item ``Limitations'' and ``future research'' sections tell you what is missing.
      \item \key{Survey papers}: a map of the field and its gaps.
      \item Replicate, then extend: new data, market, period, or a stronger baseline.
    \end{itemize}
    \column{0.46\textwidth}
    \textbf{From the world}
    \begin{itemize}
      \item A decision someone has to make: a risk limit, a price, a stock level.
      \item A new method for an old problem; an old method failing on new data.
      \item Competitions and benchmarks: M-competitions, Kaggle.
    \end{itemize}
  \end{columns}

  \begin{block}{A good question is}
    \small specific, connected to a literature, feasible with data you can get, interesting either way.
  \end{block}
\end{frame}

\begin{frame}{Finding papers and judging journals}
  \begin{columns}[T,onlytextwidth]
    \column{0.47\textwidth}
    \textbf{Search}
    \begin{itemize}
      \item \key{Google Scholar}: start here; use ``Cited by'' and ``Related articles''.
      \item \key{Scopus, Web of Science}: through the University Library (BUW), also from home.
      \item Semantic Scholar, Connected Papers: citation graphs.
      \item Preprints: arXiv, SSRN, RePEc/IDEAS.
      \item \key{Snowballing}: backwards through references, forwards through citations.
    \end{itemize}
    \column{0.5\textwidth}
    \textbf{Is this journal any good?}
    \begin{itemize}
      \item Where are the papers \emph{you} cite published? That is your shortlist.
      \item Rankings: Scimago (SJR quartiles), the Polish ministerial list of journals (points),
            the ABS Academic Journal Guide.
      \item Read the journal's \key{aims and scope} and skim two recent issues.
      \item \must{Beware predatory journals}: unsolicited invitations, fees with no real review,
            a scope that covers everything.
    \end{itemize}
  \end{columns}
\end{frame}

\begin{frame}{Journals and venues close to this seminar}
  \footnotesize
  \begin{tabular}{@{}>{\raggedright}p{0.25\textwidth}>{\raggedright\arraybackslash}p{0.71\textwidth}@{}}
    \toprule
    \textbf{Area} & \textbf{Where to read, and later where to submit} \\
    \midrule
    Forecasting & International Journal of Forecasting; Journal of Forecasting \\
    \addlinespace
    Financial econometrics, risk & Journal of Financial Econometrics; Journal of Econometrics;
      Quantitative Finance; Journal of Banking \& Finance; Insurance: Mathematics and Economics \\
    \addlinespace
    Applied ML and operations & European Journal of Operational Research; Expert Systems with
      Applications; Decision Support Systems \\
    \addlinespace
    Machine learning & NeurIPS, ICML, ICLR (conferences); Journal of Machine Learning Research;
      Transactions on Machine Learning Research \\
    \addlinespace
    Close to home & Central European Economic Journal (the journal of our faculty, open access);
      WNE Working Papers \\
    \bottomrule
  \end{tabular}

  \medskip
  \normalsize
  Pick a \key{target journal early}: it fixes format, length, audience and the standard of evidence.
\end{frame}

\begin{frame}{How to read a paper}
  \framesubtitle{Three passes, each with a decision at the end (Keshav, 2007)}
  \begin{columns}[T,onlytextwidth]
    \column{0.58\textwidth}
    \begin{enumerate}
      \item \key{Five minutes}: title, abstract, introduction, headings, conclusions.
            \muted{Relevant? Most papers stop here.}
      \item \key{An hour}: read properly, skip proofs; study figures and tables.
            \muted{Can you state the main claim and its evidence?}
      \item \key{Deep}: re-create the argument, question every assumption.
            \muted{Only for the few papers you build on.}
    \end{enumerate}
    \column{0.38\textwidth}
    \begin{block}{Keep a literature matrix}
      \small One row per paper: question, data, method, result, limitation, relation to yours.
      It grows into your literature review.
    \end{block}
    \smallskip
    \small Reference manager from day one: \key{Zotero} with Better BibTeX exports to your \texttt{.bib}.
  \end{columns}
\end{frame}

% =============================================================================
\section{Writing the article}
% =============================================================================

\begin{frame}{Anatomy of an empirical article}
  \footnotesize
  \begin{tabular}{@{}>{\raggedright\bfseries}p{0.2\textwidth}>{\raggedright}p{0.4\textwidth}>{\raggedright\arraybackslash}p{0.34\textwidth}@{}}
    \toprule
    Section & The question it answers & Typical mistake \\
    \midrule
    Introduction & What is the problem, why does it matter, what do you add? & A history of the field, no contribution \\
    Literature review & What is known, and where is the gap? & Summaries, paper by paper \\
    Data (materials) & What did you observe, and how was it prepared? & Undocumented cleaning \\
    Methods & What did you do, precisely enough to repeat it? & Textbook chapter; no evaluation protocol \\
    Results & What did you find? & Tables without interpretation \\
    Discussion & What does it mean, and what are the limits? & Claiming more than was shown \\
    Conclusions & What should the reader remember? & New material appearing here \\
    \bottomrule
  \end{tabular}

  \smallskip
  \small
  Read most of all: \key{title, abstract, figures and tables}. Many readers see nothing else.
\end{frame}

\begin{frame}{The introduction: five moves}
  \framesubtitle{It is the most-read section after the abstract. Write it to a formula (Head's ``introduction formula'').}
  \begin{enumerate}
    \item \key{Hook}: the real-world problem and why anyone should care.
    \item \key{Question}: what this paper does, in one or two sentences. Be concrete.
    \item \key{Antecedents}: the few closest papers you build on. Not the whole review.
    \item \key{Value added}: what is new here, stated as findings. Numbers are welcome:
          ``reduces the ES forecast loss by 8\% relative to GARCH''.
    \item \key{Roadmap}: how the rest of the paper is organised.
  \end{enumerate}

  \bigskip
  \begin{itemize}
    \item Test: can you state the \key{one central contribution} in a single paragraph? If not,
          the problem is the research, not the writing (Cochrane, \emph{Writing Tips for Ph.D.\ Students}).
    \item The reader should know your main result by the end of page one.
  \end{itemize}
\end{frame}

\begin{frame}{Literature review, data, methods}
  \begin{columns}[T,onlytextwidth]
    \column{0.32\textwidth}
    \textbf{Literature review}
    \begin{itemize}\small
      \item Organise by \key{idea or strand}, never paper by paper.
      \item Each strand ends with what it leaves open.
      \item The final paragraph states the gap your paper fills.
      \item Cite what you have read, at least to the second pass.
    \end{itemize}
    \column{0.32\textwidth}
    \textbf{Data}
    \begin{itemize}\small
      \item Source, period, frequency, sample size, access date.
      \item Every cleaning step and every exclusion, with counts.
      \item Descriptive statistics and one honest plot.
      \item What the data \emph{cannot} tell you.
    \end{itemize}
    \column{0.32\textwidth}
    \textbf{Methods}
    \begin{itemize}\small
      \item Your models \emph{and} the \key{baselines}, treated with equal care.
      \item The \key{evaluation protocol}: splits, windows, loss functions, tests.
      \item Notation defined once; details that interrupt go to an appendix.
      \item Enough to re-implement, with the repository as the final word.
    \end{itemize}
  \end{columns}
\end{frame}

\begin{frame}{Results, discussion, conclusions}
  \begin{columns}[T,onlytextwidth]
    \column{0.5\textwidth}
    \textbf{Results: what you found}
    \begin{itemize}
      \item Follow the order of the research questions.
      \item One message per table or figure, stated in the text.
      \item Report \key{uncertainty}: are differences statistically meaningful
            (e.g.\ Diebold--Mariano tests, model confidence sets)?
      \item \key{Robustness}: other periods, losses, hyperparameters; ablations.
      \item Report what did not work. It is a result.
    \end{itemize}
    \column{0.46\textwidth}
    \textbf{Discussion and conclusions: what it means}
    \begin{itemize}
      \item Interpret in \key{economic terms}: what changes for the risk manager, the pricing
            actuary, the planner?
      \item Compare with earlier findings: agree, disagree, why?
      \item \key{Limitations}, stated plainly, before a reviewer finds them.
      \item Conclusions: the takeaway, its practical use, what to do next. No new results.
    \end{itemize}
  \end{columns}
\end{frame}

\begin{frame}{The order you write in is not the order they read in}
  \begin{center}
  \begin{tikzpicture}[
      node distance=2.4mm,
      stage/.style={draw=semRule,fill=semTint,rounded corners=2pt,minimum height=10mm,
                    text width=15mm,align=center,font=\scriptsize},
      arr/.style={-{Stealth[length=1.8mm]},semAccent,line width=0.8pt}]
    \node[stage] (a) {\key{Figures and tables}};
    \node[stage,right=of a] (b) {\key{Data, methods}};
    \node[stage,right=of b] (c) {\key{Results}};
    \node[stage,right=of c] (d) {\key{Discussion}};
    \node[stage,right=of d] (e) {\key{Literature review}};
    \node[stage,right=of e] (f) {\key{Introduction}};
    \node[stage,right=of f] (g) {\key{Abstract, title}};
    \foreach \x/\y in {a/b,b/c,c/d,d/e,e/f,f/g} \draw[arr] (\x) -- (\y);
  \end{tikzpicture}
  \end{center}

  \medskip
  \begin{itemize}
    \item Start from the \key{exhibits}: if four figures and tables tell the story, the paper exists.
    \item Methods and data can be drafted while experiments run. Keep them current.
    \item The introduction and abstract are rewritten last, when you know what you found.
    \item Draft fast, revise slowly. Expect \key{several full revisions}; the first draft is for you,
          the later ones for the reader.
    \item One idea per paragraph, the point in its first sentence.
  \end{itemize}
\end{frame}

% =============================================================================
\section{Doing it right}
% =============================================================================

\begin{frame}{Pitfalls that invalidate results in ML for finance}
  \begin{columns}[T,onlytextwidth]
    \column{0.5\textwidth}
    \begin{itemize}
      \item \must{Look-ahead bias and leakage}: scaling or tuning that saw the test period.
      \item \must{Random splits on time series}: use rolling or expanding windows.
      \item \must{Peeking at the test set} until something works. Touch it once.
      \item \key{Backtest overfitting}: try enough strategies and one shines by chance.
    \end{itemize}
    \column{0.46\textwidth}
    \begin{itemize}
      \item \key{Weak baselines}: an untuned GARCH or a default XGBoost proves nothing.
      \item \key{Survivorship bias}: only assets or firms that still exist.
      \item \key{One lucky seed}: report variation across seeds and windows.
      \item \key{Metric mismatch}: a loss that ignores the decision.
    \end{itemize}
  \end{columns}

  \begin{block}{Worth reading before your first experiment}
    \footnotesize Kapoor \& Narayanan (2023), leakage and the reproducibility crisis, \emph{Patterns};
    Lones (2024), avoiding common ML pitfalls, \emph{Patterns};
    Bailey et al.\ (2014), backtest overfitting, \emph{Notices of the AMS};
    Bergmeir \& Benítez (2012), cross-validation for time series, \emph{Information Sciences}.
  \end{block}
\end{frame}

\begin{frame}{Working habits that make the difference}
  \begin{columns}[T,onlytextwidth]
    \column{0.48\textwidth}
    \textbf{Every day}
    \begin{itemize}
      \item A \key{research log}: date, what you tried, what happened, what you concluded.
            A markdown file in the repository is enough.
      \item Experiments driven by \key{configuration files}, results saved with the config and
            the commit hash.
      \item Write a little every week, even when ``there is nothing to write yet''.
    \end{itemize}
    \column{0.48\textwidth}
    \textbf{When stuck}
    \begin{itemize}
      \item Ask early, and well: what you wanted, what you did, what happened, what you tried.
      \item Shrink the problem: fewer series, a shorter period, one model.
      \item Reproduce a published result first. It calibrates your pipeline and your expectations.
    \end{itemize}
  \end{columns}

  \smallskip
  \textbf{Integrity is not negotiable:} your own words, every source cited, no fabricated or
  selectively reported results, AI use declared.
\end{frame}

% =============================================================================
\section{From draft to publication}
% =============================================================================

\begin{frame}{What happens after the draft}
  \begin{center}
  \begin{tikzpicture}[
      node distance=3.2mm,
      stage/.style={draw=semRule,fill=semTint,rounded corners=2pt,minimum height=11mm,
                    text width=17.5mm,align=center,font=\scriptsize},
      arr/.style={-{Stealth[length=2mm]},semAccent,line width=0.9pt}]
    \node[stage] (a) {\key{Internal review}\\ supervisor, peers};
    \node[stage,right=of a] (b) {\key{Preprint}\\ arXiv, SSRN};
    \node[stage,right=of b] (c) {\key{Submission}\\ one journal at a time};
    \node[stage,right=of c] (d) {\key{Editor}\\ desk reject or review};
    \node[stage,right=of d] (e) {\key{Referee reports}\\ revise and resubmit};
    \node[stage,right=of e] (f) {\key{Accepted}};
    \foreach \x/\y in {a/b,b/c,c/d,d/e,e/f} \draw[arr] (\x) -- (\y);
  \end{tikzpicture}
  \end{center}

  \medskip
  \begin{itemize}
    \item Peer review takes \key{months per round}. Rejection is the most common outcome at good
          journals; you revise and submit elsewhere.
    \item A revision comes with a \key{response letter}: every comment quoted and answered.
    \item \key{Your thesis does not wait for this.} Publication is not required before the defence.
    \item Conferences and seminars are where you test the paper before referees do.
  \end{itemize}
\end{frame}

\begin{frame}{First-timer's checklist}
  \begin{columns}[T,onlytextwidth]
    \column{0.48\textwidth}
    \textbf{This month}
    \begin{enumerate}
      \item Install Zotero; start your literature matrix.
      \item Skim the last two years of one target journal.
      \item Read three papers to pass two, one to pass three.
      \item Start the research log in your repository.
      \item Draft the two-page concept note.
    \end{enumerate}
    \column{0.48\textwidth}
    \textbf{Before the first experiment}
    \begin{enumerate}
      \item Write down the evaluation protocol; freeze the test period.
      \item Implement and tune the baselines first.
      \item Reproduce one published number.
      \item Sketch the tables and figures you expect to need.
    \end{enumerate}
  \end{columns}

  \smallskip
  \muted{\small More in the repository: \texttt{docs/research-onboarding.md} (reading list, data sources, tools, glossary).}
\end{frame}

\end{document}
